% !TeX root = ../main.tex

\begin{resume}

  \section*{个人简历}

  林俊峰，1999 年 1 月 23 日出生于四川。

  2017 年 9 月考入清华大学计算机科学与技术专业，2021 年 6 月本科毕业并获得工学学士学位。

  2021 年 9 月免试进入清华大学类脑计算中心攻读博士学位至今。

  2025 年 2 月至 2025 年 8 月，以访问学者身份赴伊利诺伊大学-厄巴纳香槟分校进行学术交流。


  \section*{在学期间完成的相关学术成果}

  \subsection*{学术论文}

  \begin{achievements}
    \item LIN J, QU H, MA S, et al. SongC: A Compiler for Hybrid Near-Memory and In-Memory Many-Core Architecture[J]. IEEE Transactions on Computers, 2023. (CCF-A类期刊，第一作者)
    \item LIN J, LIU Z, et al. WeiPipe: Weight Pipeline Parallelism for Communication-Effective Long-Context Large Model Training[C]. ACM SIGPLAN Symposium on Principles and Practice of Parallel Programming (PPoPP), 2025. (CCF-A类会议，口头报告)
    \item YU J, LIN J, et al. VoltanaLLM: Feedback-Driven Frequency Control and State-Space Routing for Energy-Efficient LLM Serving[C]. International Symposium on Computer Architecture (ISC), 2026. (CCF-B类会议)
    \item QU H, ZHANG W, LIN J, et al. CATFlow: Data-Monistic Parallelism Plan Discovery for Distributed Training[C]. USENIX Symposium on Operating Systems Design and Implementation (SOSP), 2026. (CCF-A类会议，在审)
  \end{achievements}


  \subsection*{专利}

  \begin{achievements}
    \item 施路平，林俊峰，等. 权重流水并行的神经网络分布式训练方法和装置[P]: 北京市，20231 0160001.1. 2024-07-19.
    \item 施路平，林俊峰，等. 任务处理的方法和装置、众核系统、计算机可读介质[P]: 北京市， 2023 1 0160001.1. 2023-02-21.
    \item 施路平，林俊峰，等. 神经网络任务的代码生成方法和装置[P]: 北京市, ZL 2022 1 0705825.8. 2022-06-21.
  \end{achievements}


\end{resume}